\section*{Problems}

\subsection{Problem Statements}

\subsubsection*{1. Fundamental Level (Simple / Direct)}

\begin{problema}
	\label{prob:2.6.1}
	A finite population consists of $N = 6$ distinct elements $\{a, b, c, d, e, f\}$. A random sample of size $n = 3$ is to be drawn.
	\begin{enumerate}
		\item Determine the total number of possible samples if the draw is performed \textbf{with replacement}.
		\item Determine the total number of possible samples if the draw is performed \textbf{without replacement} and without regard to order (unordered samples).
		\item Formally explain why, as the population size $N$ grows toward infinity, the probabilistic difference between sampling with and without replacement tends to zero.
	\end{enumerate}
\end{problema}

\begin{problema}
	\label{prob:2.6.2}
	Let $X_1, X_2, \ldots, X_n$ be a simple random sample (independent and identically distributed random variables) drawn from a population with theoretical mean $\mu$ and variance $\sigma^2 < \infty$. Prove algebraically, applying the properties of the mathematical expectation operator, that the sample mean $\bar{X} = \frac{1}{n} \sum_{i=1}^n X_i$ is an \textbf{unbiased estimator} of the population parameter $\mu$, that is, that $E(\bar{X}) = \mu$.
\end{problema}

\begin{problema}
	\label{prob:2.6.3}
	The assembly time of a mechanical component follows a distribution with population mean $\mu = 45$ minutes and standard deviation $\sigma = 12$ minutes. If a random sample of size $n$ is taken, compute the standard error of the sample mean, denoted by $\text{SE}(\bar{X}) = \sigma_{\bar{X}}$, for each of the following sample sizes and comment on the rate at which the error decreases:
	\begin{enumerate}
		\item $n = 9$
		\item $n = 36$
		\item $n = 144$
		\item $n = 576$
	\end{enumerate}
\end{problema}


\subsubsection*{2. Operational Level (Intermediate / Applied)}

\begin{problema}
	\label{prob:2.6.4}
	A government audit evaluates the monthly income of a trade guild that has exactly $N = 400$ registered companies. The population standard deviation is historically known to be $\sigma = \$15,000$ MXN. An audit team selects a random sample of $n = 64$ companies.
	\begin{enumerate}
		\item Compute the variance of the sample mean $\text{Var}(\bar{X})$ under the assumption of sampling \textbf{with replacement} (infinite population or multinomial sampling).
		\item Compute the variance of the sample mean $\text{Var}(\bar{X})$ under sampling \textbf{without replacement}, applying the finite population correction (FPC) factor given by $\frac{N-n}{N-1}$.
		\item What percentage reduction in the sample variance is achieved thanks to the finite population correction in this audit scenario?
	\end{enumerate}
\end{problema}

\begin{problema}
	\label{prob:2.6.5}
	A fair 6-sided die is rolled $n = 36$ times independently. For a single roll, the population mean is $\mu = 3.5$ and the variance is $\sigma^2 = \frac{35}{12} \approx 2.9167$. Let $\bar{X}_{36}$ be the arithmetic average of the 36 rolls.
	Using the Central Limit Theorem (CLT), compute the approximate probability that the sample average obtained lies strictly between $3.1$ and $3.9$, that is, $P(3.1 < \bar{X}_{36} < 3.9)$.
\end{problema}

\begin{problema}
	\label{prob:2.6.6}
	A mineral water bottling plant wants to calibrate its production line. The amount poured into each bottle is a random variable with known standard deviation $\sigma = 8$ milliliters.
	What is the minimum sample size $n$ that a quality engineer must inspect to have at least $95\%$ probabilistic certainty that the sample mean $\bar{X}$ does not differ from the true population filling mean $\mu$ by more than $1.5$ milliliters?
\end{problema}


\subsubsection*{3. Analytical Level (Advanced / Proofs)}

\begin{problema}
	\label{prob:2.6.7}
	Prove the \textbf{Weak Law of Large Numbers (WLLN)} using Chebyshev's inequality. Specifically, prove that if $X_1, X_2, \ldots, X_n$ are i.i.d. random variables with mean $\mu$ and variance $\sigma^2 < \infty$, then for any arbitrary tolerance $\epsilon > 0$, the probability that the sample mean $\bar{X}_n$ is $\epsilon$ or more away from $\mu$ decreases to zero as $n \to \infty$:
	\begin{align}
		\lim_{n \to \infty} P(|\bar{X}_n - \mu| \ge \epsilon) = 0
	\end{align}
\end{problema}

\begin{problema}
	\label{prob:2.6.8}
	Prove algebraically why the sample variance estimator requires dividing by $n-1$ (and not by $n$) in order to be unbiased. Let $S^2 = \frac{1}{n-1} \sum_{i=1}^n (X_i - \bar{X})^2$ be the sample variance of a SRS drawn from a population with mean $\mu$ and variance $\sigma^2$. Prove step by step that $E(S^2) = \sigma^2$.
	\emph{Hint:} Add and subtract $\mu$ inside the squared term $(X_i - \bar{X})^2 = [(X_i - \mu) - (\bar{X} - \mu)]^2$.
\end{problema}


\subsubsection*{4. Challenging Level (Expert / Paradoxes)}

\begin{problema}
	\label{prob:2.6.9}
	The Central Limit Theorem postulates that the sum of independent variables will converge to the normal distribution as long as the second-order moment ($\sigma^2$) is finite. Analyze what happens when samples are drawn from a population with extremely heavy tails: the standard Cauchy distribution, whose density is $f(x) = \frac{1}{\pi(1+x^2)}$ for $x \in (-\infty, \infty)$.
	\begin{enumerate}
		\item Formally prove or argue why the integral defining the theoretical expectation $E(X)$ does not converge absolutely for the Cauchy distribution.
		\item If a random sample of size $n$ is drawn from a Cauchy population, what is the density function or behavior of the sample mean $\bar{X}_n$? Does convergence toward the normal occur as $n \to \infty$?
	\end{enumerate}
\end{problema}

\begin{problema}
	\label{prob:2.6.10}
	In computational practice, the Central Limit Theorem provides an asymptotic approximation. To evaluate the exact speed and quality of this convergence for finite samples $n$, the \textbf{Berry-Esseen Theorem} is used.
	Let $X_1, X_2, \ldots, X_n$ be an i.i.d. sample with $E(X_i) = \mu$, $\text{Var}(X_i) = \sigma^2 > 0$, and finite absolute third central moment $\rho_3 = E[|X_i - \mu|^3] < \infty$. The Berry-Esseen bound establishes that the maximum absolute approximation error between the standardized empirical distribution function $F_n(z)$ and the standard normal distribution $\Phi(z)$ satisfies:
	\begin{align}
		\sup_{z \in \mathbb{R}} |F_n(z) - \Phi(z)| \le \frac{C \cdot \rho_3}{\sigma^3 \sqrt{n}}
	\end{align}
	where $C < 0.4784$ is the universal Berry-Esseen constant.
	If an exponential population with parameter $\lambda = 1$ ($\mu=1, \sigma=1, \rho_3 \approx 3.6$) is sampled with $n = 100$, compute the upper bound of the absolute probability error guaranteed by the Berry-Esseen Theorem and interpret its significance for engineering decision-making.
\end{problema}


\subsection{Detailed Solutions}

\subsubsection*{1. Fundamental Level (Simple / Direct)}

\begin{solproblema}
	\begin{enumerate}
		\item \textbf{Sampling with replacement:} In each of the $n = 3$ draws there are exactly $N = 6$ available options, since the drawn element is returned to the population before the next selection. By the multiplicative counting principle, the total number of possible ordered samples is:
		\begin{align}
			N^n = 6^3 = 216 \text{ samples.}
		\end{align}

		\item \textbf{Sampling without replacement (unordered):} Since each element can only appear once per sample and the order in which elements are drawn does not alter the final set, the total number of possible subsets is given by the binomial coefficient:
		\begin{align}
			\binom{N}{n} = \binom{6}{3} = \frac{6!}{3!(6-3)!} = \frac{6 \times 5 \times 4}{3 \times 2 \times 1} = 20 \text{ samples.}
		\end{align}

		\item \textbf{Asymptotic equivalence as $N \to \infty$:} When drawing a fixed sample of size $n$ without replacement, the probability of selecting a specific element on the first draw is $1/N$, on the second is $1/(N-1)$, and on the $k$-th is $1/(N-k+1)$. As $N \to \infty$ with $n$ held constant, the ratios $\frac{N-k}{N} \to 1$. Therefore, population depletion becomes negligible and sampling without replacement converges in distribution to sampling with replacement (independent).
	\end{enumerate}
\end{solproblema}

\begin{solproblema}
	By hypothesis, $X_1, \ldots, X_n$ are independent and identically distributed, so that for every $i \in \{1, \ldots, n\}$ we have $E(X_i) = \mu$. Applying the linearity of the expectation operator (the expectation of a sum is the sum of expectations, and multiplicative constants factor out of the expectation):
	\begin{align}
		E(\bar{X}) &= E\left( \frac{1}{n} \sum_{i=1}^n X_i \right) \\
		&= \frac{1}{n} \sum_{i=1}^n E(X_i) \\
		&= \frac{1}{n} \sum_{i=1}^n \mu \\
		&= \frac{1}{n} (n \cdot \mu) = \mu.
	\end{align}
	Since $E(\bar{X}) = \mu$, it is algebraically demonstrated that the sample average is an unbiased estimator of the population mean.
\end{solproblema}

\begin{solproblema}
	The standard error of the mean ($\sigma_{\bar{X}}$) measures the dispersion of sample averages around $\mu$ and is given by $\sigma_{\bar{X}} = \frac{\sigma}{\sqrt{n}}$. With $\sigma = 12$:
	\begin{enumerate}
		\item For $n = 9$: $\sigma_{\bar{X}} = \frac{12}{\sqrt{9}} = \frac{12}{3} = 4.00 \text{ min.}$
		\item For $n = 36$: $\sigma_{\bar{X}} = \frac{12}{\sqrt{36}} = \frac{12}{6} = 2.00 \text{ min.}$
		\item For $n = 144$: $\sigma_{\bar{X}} = \frac{12}{\sqrt{144}} = \frac{12}{12} = 1.00 \text{ min.}$
		\item For $n = 576$: $\sigma_{\bar{X}} = \frac{12}{\sqrt{576}} = \frac{12}{24} = 0.50 \text{ min.}$
	\end{enumerate}
	\textbf{Comment on the rate of reduction:} The dispersion decreases strictly in inverse proportion to the square root of the sample size ($\sqrt{n}$). To halve the standard error (for example, from 4.0 to 2.0, or from 2.0 to 1.0), the inspected sample size must be quadrupled ($9 \to 36 \to 144$).
\end{solproblema}


\subsubsection*{2. Operational Level (Intermediate / Applied)}

\begin{solproblema}
	Known data: $N = 400$, $\sigma = 15,000$ MXN, and $n = 64$.
	\begin{enumerate}
		\item \textbf{Variance under sampling with replacement (or infinite population):}
		\begin{align}
			\text{Var}_{W/R}(\bar{X}) = \frac{\sigma^2}{n} = \frac{(15,000)^2}{64} = \frac{225,000,000}{64} = 3,515,625 \text{ MXN}^2.
		\end{align}
		The standard error with replacement would be $\sigma_{\bar{X}} = \sqrt{3,515,625} = \$1,875 \text{ MXN}$.

		\item \textbf{Variance under sampling without replacement (finite population):}
		We first compute the finite population correction (FPC) factor:
		\begin{align}
			\text{FPC} = \frac{N - n}{N - 1} = \frac{400 - 64}{400 - 1} = \frac{336}{399} \approx 0.842105.
		\end{align}
		We multiply the sample variance by the FPC:
		\begin{align}
			\text{Var}_{W/o R}(\bar{X}) = \frac{\sigma^2}{n} \left( \frac{N-n}{N-1} \right) = 3,515,625 \times 0.842105 \approx 2,960,156.25 \text{ MXN}^2.
		\end{align}
		The actual standard error without replacement is $\sigma_{\bar{X}, W/o R} = \sqrt{2,960,156.25} \approx \$1,720.51 \text{ MXN}$.

		\item \textbf{Percentage reduction in variance:}
		The relative reduction in variance is exactly equal to $1 - \text{FPC}$:
		\begin{align}
			\text{Reduction} = 1 - \frac{336}{399} = \frac{63}{399} \approx 0.15789 \text{ (15.79\%).}
		\end{align}
		By sampling $16\%$ of the population ($64/400$), the precision gain from not repeating companies reduces the variance by $15.79\%$.
	\end{enumerate}
\end{solproblema}

\begin{solproblema}
	For the fair die: $\mu = 3.5$, $\sigma^2 = \frac{35}{12} \approx 2.9167$, $n = 36$.
	The standard error of the average of 36 rolls is:
	\begin{align}
		\sigma_{\bar{X}} = \frac{\sigma}{\sqrt{n}} = \sqrt{\frac{35/12}{36}} = \sqrt{\frac{35}{432}} \approx \sqrt{0.0810185} \approx 0.284637.
	\end{align}
	We want to compute $P(3.1 < \bar{X}_{36} < 3.9)$. We standardize the interval limits using the variable $Z = \frac{\bar{X} - \mu}{\sigma_{\bar{X}}}$:
	\begin{align}
		z_1 &= \frac{3.1 - 3.5}{0.284637} = \frac{-0.40}{0.284637} \approx -1.4053 \\
		z_2 &= \frac{3.9 - 3.5}{0.284637} = \frac{+0.40}{0.284637} \approx +1.4053
	\end{align}
	By the CLT, $\bar{X}_{36}$ approaches a normal distribution, so:
	\begin{align}
		P(3.1 < \bar{X}_{36} < 3.9) &\approx P(-1.405 < Z < +1.405) \\
		&= \Phi(1.405) - \Phi(-1.405) = 2\Phi(1.405) - 1.
	\end{align}
	Using standard normal tables, $\Phi(1.405) \approx 0.9200$. Therefore:
	\begin{align}
		P(3.1 < \bar{X}_{36} < 3.9) \approx 2(0.9200) - 1 = 0.8400 \text{ (84.00\%).}
	\end{align}
\end{solproblema}

\begin{solproblema}
	Data: $\sigma = 8$ ml, maximum tolerated error margin $E = 1.5$ ml, $95\%$ confidence level ($\alpha = 0.05$).
	For a $95\%$ probability, the critical standard normal value that accumulates $0.975$ in the upper tail is $z_{\alpha/2} = 1.96$.
	The margin of error for estimating the mean is:
	\begin{align}
		E = z_{\alpha/2} \cdot \frac{\sigma}{\sqrt{n}}
	\end{align}
	Solving for the minimum sample size $n$:
	\begin{align}
		\sqrt{n} &= \frac{z_{\alpha/2} \cdot \sigma}{E} \\
		n &= \left( \frac{z_{\alpha/2} \cdot \sigma}{E} \right)^2 = \left( \frac{1.96 \times 8}{1.5} \right)^2 = \left( \frac{15.68}{1.5} \right)^2 \approx (10.4533)^2 \approx 109.27.
	\end{align}
	Since $n$ must be an integer and at least $95\%$ certainty must be guaranteed, we always round up to the next integer:
	\begin{align}
		n_{\text{min}} = 110 \text{ bottles.}
	\end{align}
\end{solproblema}


\subsubsection*{3. Analytical Level (Advanced / Proofs)}

\begin{solproblema}
	Let $X_1, \ldots, X_n$ be a sequence of i.i.d. random variables with $E(X_i) = \mu$ and $\text{Var}(X_i) = \sigma^2$.
	We know from the properties of expectation and variance that for the sample mean $\bar{X}_n = \frac{1}{n}\sum_{i=1}^n X_i$:
	\begin{align}
		E(\bar{X}_n) &= \mu, \\
		\text{Var}(\bar{X}_n) &= \frac{\sigma^2}{n}.
	\end{align}
	\textbf{Chebyshev's Inequality} establishes that for any random variable $Y$ with mean $\mu_Y$ and variance $\sigma_Y^2$, and for every constant $k > 0$ or tolerance $\epsilon > 0$:
	\begin{align}
		P(|Y - \mu_Y| \ge \epsilon) \le \frac{\text{Var}(Y)}{\epsilon^2}.
	\end{align}
	Substituting $Y = \bar{X}_n$, $\mu_Y = \mu$, and $\text{Var}(Y) = \frac{\sigma^2}{n}$ into the inequality:
	\begin{align}
		0 \le P(|\bar{X}_n - \mu| \ge \epsilon) \le \frac{\sigma^2/n}{\epsilon^2} = \frac{\sigma^2}{n \epsilon^2}.
	\end{align}
	Taking the limit on both ends as the sample size $n \to \infty$ (treating $\sigma^2$ and $\epsilon > 0$ as fixed constants):
	\begin{align}
		\lim_{n \to \infty} \frac{\sigma^2}{n \epsilon^2} = 0.
	\end{align}
	By the Squeeze Theorem for algebraic limits, we rigorously conclude that:
	\begin{align}
		\lim_{n \to \infty} P(|\bar{X}_n - \mu| \ge \epsilon) = 0,
	\end{align}
	which proves that the sample mean $\bar{X}_n$ converges in probability to the true mean $\mu$ ($\bar{X}_n \xrightarrow{p} \mu$).
\end{solproblema}

\begin{solproblema}
	We want to find $E(S^2)$ where $S^2 = \frac{1}{n-1} \sum_{i=1}^n (X_i - \bar{X})^2$.
	First, we analyze the sum of squared deviations from the sample mean $\sum_{i=1}^n (X_i - \bar{X})^2$. We add and subtract the true population mean $\mu$ inside the parenthesis:
	\begin{align}
		\sum_{i=1}^n (X_i - \bar{X})^2 &= \sum_{i=1}^n [(X_i - \mu) - (\bar{X} - \mu)]^2 \\
		&= \sum_{i=1}^n \left[ (X_i - \mu)^2 - 2(X_i - \mu)(\bar{X} - \mu) + (\bar{X} - \mu)^2 \right] \\
		&= \sum_{i=1}^n (X_i - \mu)^2 - 2(\bar{X} - \mu) \sum_{i=1}^n (X_i - \mu) + n(\bar{X} - \mu)^2.
	\end{align}
	Note that the linear sum with respect to $\mu$ can be simplified: $\sum_{i=1}^n (X_i - \mu) = \sum_{i=1}^n X_i - n\mu = n\bar{X} - n\mu = n(\bar{X} - \mu)$.
	Substituting this into the middle term:
	\begin{align}
		\sum_{i=1}^n (X_i - \bar{X})^2 &= \sum_{i=1}^n (X_i - \mu)^2 - 2(\bar{X} - \mu)[n(\bar{X} - \mu)] + n(\bar{X} - \mu)^2 \\
		&= \sum_{i=1}^n (X_i - \mu)^2 - 2n(\bar{X} - \mu)^2 + n(\bar{X} - \mu)^2 \\
		&= \sum_{i=1}^n (X_i - \mu)^2 - n(\bar{X} - \mu)^2.
	\end{align}
	Now we apply the expectation operator $E(\cdot)$ to both sides of this identity:
	\begin{align}
		E\left[ \sum_{i=1}^n (X_i - \bar{X})^2 \right] &= \sum_{i=1}^n E[(X_i - \mu)^2] - n \cdot E[(\bar{X} - \mu)^2].
	\end{align}
	By definition of the theoretical population variance and the variance of the sample mean:
	\begin{itemize}
		\item $E[(X_i - \mu)^2] = \text{Var}(X_i) = \sigma^2$ for each of the $n$ terms.
		\item $E[(\bar{X} - \mu)^2] = \text{Var}(\bar{X}) = \frac{\sigma^2}{n}$.
	\end{itemize}
	Substituting these exact variances:
	\begin{align}
		E\left[ \sum_{i=1}^n (X_i - \bar{X})^2 \right] &= \sum_{i=1}^n \sigma^2 - n \left( \frac{\sigma^2}{n} \right) \\
		&= n\sigma^2 - \sigma^2 = (n - 1)\sigma^2.
	\end{align}
	Finally, dividing by the normalization constant $(n-1)$:
	\begin{align}
		E(S^2) = E\left[ \frac{1}{n-1} \sum_{i=1}^n (X_i - \bar{X})^2 \right] = \frac{1}{n-1} E\left[ \sum_{i=1}^n (X_i - \bar{X})^2 \right] = \frac{(n-1)\sigma^2}{n-1} = \sigma^2.
	\end{align}
	\textbf{Conclusion:} If one divided by $n$ instead, the estimator would have expectation $\frac{n-1}{n}\sigma^2$, systematically underestimating the population variance. Dividing by the $n-1$ degrees of freedom exactly corrects this bias.
\end{solproblema}


\subsubsection*{4. Challenging Level (Expert / Paradoxes)}

\begin{solproblema}
	\begin{enumerate}
		\item \textbf{Nonexistence of moments (Non-convergent expectation):}
		The expectation of a continuous random variable with density $f(x) = \frac{1}{\pi(1+x^2)}$ is given by:
		\begin{align}
			E(X) = \int_{-\infty}^{\infty} x \cdot \frac{1}{\pi(1+x^2)} dx = \frac{1}{\pi} \int_{-\infty}^{\infty} \frac{x}{1+x^2} dx.
		\end{align}
		For an improper integral over both extremes to converge in the Lebesgue sense, the integrals of the positive part ($x \ge 0$) and negative part ($x < 0$) must be absolutely finite separately:
		\begin{align}
			\int_0^{\infty} \frac{x}{1+x^2} dx = \lim_{b \to \infty} \left[ \frac{1}{2} \ln(1+x^2) \right]_0^b = \lim_{b \to \infty} \frac{1}{2} \ln(1+b^2) = \infty.
		\end{align}
		Since the integral over $[0, \infty)$ diverges to $+\infty$ (and analogously over $(-\infty, 0]$ diverges to $-\infty$), the total improper integral is an indeterminate form of type $\infty - \infty$. Therefore, the first-order moment $E(X)$ (and even more so the second-order moment $\sigma^2$) \textbf{does not exist}.

		\item \textbf{Distribution of the sample mean for Cauchy:}
		Let $\bar{X}_n = \frac{1}{n} \sum_{i=1}^n X_i$ where $X_i \sim \text{Cauchy}(0,1)$ i.i.d. Using characteristic functions (Fourier transform), the characteristic function of a standard Cauchy variable is $\varphi_X(t) = E(e^{itX}) = e^{-|t|}$.
		The characteristic function of the sample average $\bar{X}_n$ is:
		\begin{align}
			\varphi_{\bar{X}_n}(t) &= E\left[ \exp\left( it \frac{1}{n} \sum_{k=1}^n X_k \right) \right] = \prod_{k=1}^n \varphi_X\left( \frac{t}{n} \right) \\
			&= \prod_{k=1}^n \exp\left( -\left| \frac{t}{n} \right| \right) = \exp\left( -n \frac{|t|}{n} \right) = e^{-|t|}.
		\end{align}
		Surprisingly, the characteristic function $\varphi_{\bar{X}_n}(t) = e^{-|t|}$ is identical to the characteristic function of a single Cauchy observation!
		By the uniqueness theorem of Fourier inversion, the distribution of the average of $n$ Cauchy variables is exactly the same original Cauchy distribution:
		\begin{align}
			\bar{X}_n \sim \text{Cauchy}(0,1).
		\end{align}
		\textbf{Statistical interpretation:} In Cauchy populations (or Pareto with extreme tails $\alpha \le 1$), \textbf{no convergence toward the normal occurs} and there is no variance reduction. Taking the average of a million Cauchy data points has exactly the same dispersion, probability of extreme values, and estimation error as taking a single random data point. The Central Limit Theorem breaks down completely when the hypothesis $\sigma^2 < \infty$ is violated.
	\end{enumerate}
\end{solproblema}

\begin{solproblema}
	For a variable $X \sim \text{Exp}(\lambda=1)$: the density is $f(x) = e^{-x}$ for $x \ge 0$.
	\begin{itemize}
		\item Population mean: $\mu = E(X) = 1$.
		\item Standard deviation: $\sigma = \sqrt{\text{Var}(X)} = 1$.
		\item Absolute third central moment: $\rho_3 = E[|X - \mu|^3] = \int_0^{\infty} |x - 1|^3 e^{-x} dx \approx 3.60$.
		\item Universal Berry-Esseen constant: $C < 0.4784$.
	\end{itemize}
	Evaluating the Berry-Esseen bound for a sample of size $n = 100$:
	\begin{align}
		\sup_{z \in \mathbb{R}} |F_{100}(z) - \Phi(z)| \le \frac{C \cdot \rho_3}{\sigma^3 \sqrt{n}} = \frac{0.4784 \times 3.60}{1^3 \times \sqrt{100}} = \frac{1.72224}{10} \approx 0.1722.
	\end{align}
	\textbf{Interpretation for engineering and data science:}
	The bound rigorously guarantees that when using the standard normal table $\Phi(z)$ to approximate probabilities of the sample mean with $n=100$ from a population as skewed as the exponential, the maximum possible error at any point on the curve is \textbf{$17.22\%$} (or less than $0.1722$ in absolute probability).
	This shows why the classical "$n \ge 30$ rule of thumb" is insufficient for strongly skewed populations: to reduce the maximum asymptotic CLT approximation error to below $1\%$ ($0.01$), an exponential distribution would require a sample size of at least $n \ge \left(\frac{1.72224}{0.01}\right)^2 \approx 29,661$ observations.
\end{solproblema}
